\section{Experiments}
\label{sec:experiments}
\subsection{Data and Regional Evaluation}
We use 1,713 paired Sentinel-1 locations from the Haiti landslide-change dataset \cite{Bralet2024DataPort}. Each orbit provides pre/post VV and VH plus an orbit identifier in a $128\times128$ patch. The Student uses ConvNeXt-Tiny/FPN, Adam with learning rate $5\times10^{-5}$, and batch size 16. Training includes all labelled pixels without geometry input or geometry weighting. One inference output is scored on four masks: TargetValid ($g_t=0$), Overall (all pixels), BothValid ($g_t=g_c=0$), and $G_{10}$ ($g_t>0,g_c=0$), where zero is valid geometry. AUPRC uses scores; IoU/F1 use threshold 0.5.

\subsection{Mapping Comparison}
Table~\ref{tab:primary} compares five representative baselines with the current Student. Ours improves TargetValid AUPRC/IoU over S0 and Multi2Uni and retains gains in Overall, BothValid, and $G_{10}$. External rows are target-input adaptations. Archived DIS2 retains its recorded Overall 87.33\% AUPRC and 61.77\% IoU; original probabilities and mask/sample lineage are unavailable, so its regional cells remain blank. Its historical TGD scores are not merged with current $G_{10}$.
\input{sections/current_primary_table}

\subsection{Qualitative Results}
Fig.~\ref{fig:qualitative} provides native-patch examples; Fig.~\ref{fig:g10_zoom} retains predictions on distorted pixels. Scene/crop selection uses annotation support, independently of comparative performance.

\begin{figure*}[!t]
\centering
\includegraphics[width=\textwidth]{figures/fig_qualitative_current_raw.pdf}
\caption{Scenes 001061 DESC, 000697 ASC, and 000925 ASC: target/counter VV, target geometry, annotation, and unmasked predictions. Green: TP; red: FP; blue: FN; dark: TN. Both orbits share one pre/post VV display range per scene; counter SAR is explanatory only.}
\label{fig:qualitative}
\end{figure*}

\begin{figure*}[!t]
\centering
\includegraphics[width=.78\textwidth]{figures/fig_g10_raw_zoom.pdf}
\caption{Unmasked $G_{10}$ detail, 000925 ASC, crop $x\in[24,88)$, $y\in[8,72)$. Yellow contours and the mask identify target-distorted/counter-valid pixels. The same TP/FP/FN colors apply; no prediction is replaced by gray.}
\label{fig:g10_zoom}
\end{figure*}

\subsection{Current Search Controls}
Table~\ref{tab:search_intervention} uses frozen current checkpoints. Center sets offsets to zero; random draws bounded offsets while retaining the anchored center and learned weights; no-search feeds projected P2 to the trained decoder. Counter shuffle globally deranges sample locations. Frozen CG0 self-context differs from independently trained SG0. Counter mismatch and random sampling reduce accuracy, but Student center sampling is close to learned sampling. These are inference sensitivities, not proof that a retrained fixed sampler would be inferior.

\input{sections/current_search_table}

\subsection{Policy and Shared-Memory Evidence}
Alignment diagnostics use the current CG0 and S2-2Stage checkpoints. Over fully TargetValid query blocks, candidate-weight JS is 0.0219 for Student versus CG0 and 0.0782 for shuffled versus CG0. Shared-memory cosine distance to the Student retrieval is 0.0066 for paired CG0 and 0.0595 for shuffled context. These statistics weight every eligible query equally. Candidate-weight JS is a diagnostic over twelve weights, distinct from the spatially splatted policy distribution used in the training loss.


\begin{figure*}[!t]
\centering
\includegraphics[width=.90\textwidth]{figures/fig_search_current_compact.pdf}
\caption{Current query $(57.5,49.5)$, 000925 ASC. Teacher triangles and Student circles show all four samples at each FPN level without jitter; red crosses mark the reference. Bottom: joint candidate weights and scale mass. The three learned points coincide in both checkpoints; parameter inspection confirms identical offset heads. Effective support is one learned location plus the fixed center per level.}
\label{fig:search_samples}
\end{figure*}

\input{sections/current_loss_table}
Table~\ref{tab:loss_parts} reports completed seed-42 component-removal controls with the current warm-up, CG0 guidance, and stop-gradient schedule. Each surviving component retains its canonical coefficient. Overall AUPRC is 88.69\% for segmentation alone, 88.96\% with P, 88.68\% with E, and 89.16\% for the frozen combined model. E alone does not improve this matched control. The new segmentation-only control is stronger than the archived S0; these matched results, rather than historical loss sweeps, delimit component attribution. One seed does not establish statistical superiority.


